\section{Geometry, classical solutions, and the Pompeiu conclusion}
\label{inf:sec-geometry}

\begin{lemma}[Analyticity and full-dimensional strict convexity]
\label{inf:geometry-convexity}
Let $f\in\Ec_1$ be real with
$\nrm f_\infty+\nrm{\nabla_S^2f}_{\mathrm{op},\infty}<1/4$.
The radial graph $\{(1+f(\omega))\omega:\omega\in S^{2p-1}\}$
is a closed analytic hypersurface bounding a bounded strictly
convex domain. It is invariant under $O(a)\times O(b)$.
The conclusion concerns all ambient directions, including the
block-orbit directions and the axes.
\end{lemma}
\begin{proof}
By Lemma~\ref{inf:space-embedding}, $f$ is holomorphic in a
complex neighbourhood of $x\in[0,1]$. On the real sphere
$x=|X|^2$ is a polynomial. The radial parametrization is therefore
analytic, including at $x=0,1$. Its radius is positive and its
radial derivative is nonsingular, so it is an embedded closed
hypersurface. Its symmetry follows from dependence on $x$ alone.

Put $r_f=1+f$ and let $g_S$ be the sphere metric.
Differentiating the radial parametrization and its outward normal
gives the second fundamental form with positive denominator and
numerator
\begin{equation}\label{inf:curvature-numerator}
 r_f^2g_S+2\,df\otimes df-r_f\nabla_S^2f.
\end{equation}
Since $3/4<r_f<5/4$, its quadratic form on every unit tangent
vector is at least $9/16-5/16=1/4$.
This proves positivity in the full tangent space.
To make the block directions explicit, at $0<x<1$ a unit
$X$-orbit tangent vector has
$\nabla_S^2x(v,v)=2(1-x)$ and $dx(v)=0$; a unit $Y$-orbit
tangent vector has $\nabla_S^2x(v,v)=-2x$.
Their numerators in \eqref{inf:curvature-numerator} are
$r_f^2-2r_f(1-x)f'$ and $r_f^2+2r_fx f'$.
For the unit meridian direction the numerator is
\[
 r_f^2+8x(1-x)(f')^2
   -r_f\{4x(1-x)f''+2(1-2x)f'\}.
\]
The coordinate-free positive bound already controls all three
expressions and their limits at the axes.
The global convexity theorem for compact hypersurfaces with
positive second fundamental form (Hadamard's theorem, in the form
given by \cite{inf:doCarmoLima}; see also \cite{inf:sacksteder})
now makes the enclosed radial domain
strictly convex. Its hypotheses hold: the hypersurface is smooth,
compact, connected and embedded, and every principal curvature
is positive. The convex boundary contains no line segment, since
an interior point of such a segment would have zero normal
curvature in its direction. Thus the convexity is strict.
Boundedness follows from its bounded radius.
\end{proof}

\begin{lemma}[The non-ball witness survives the correction]
\label{inf:geometry-nonball}
At the zero supplied by Corollary~\ref{inf:Newton-zero}, the
normalized shape $f_*$ satisfies the hypotheses of
Lemma~\ref{inf:geometry-convexity} for large $p$, and its domain
is not a ball.
\end{lemma}
\begin{proof}
The fixed centre has normalized $C^2$ size at most
$CC_0\tau p^{-2}$ by \eqref{inf:shape-C2},
$\nrm{f_N}_{\Ec_1}\le\delta$, and \eqref{inf:uniform-delta}.
The correction has normalized $C^2$ size at most
$C\tau p^{-L_{\mathrm{Newt}}}$ by \eqref{inf:shape-C2} and
\eqref{inf:Newton-correction}. Their sum is below $1/4$ for
large $p$, uniformly for $\tau\le20$.

Consider the physical monic seed functional
$\mathfrak W_p(f)=R[\mathsf H_2]f=Rf_2/\mathsf H_2(1)$.
Its dual norm on $\Ec_1$ is at most $Cp$, since $R=O(p)$ and
$\mathsf H_2(1)=(\bb)_2/(p+1)_2$ is bounded away from zero uniformly
in the split. At the centre,
\[
 \mathfrak W_p(f_N)
 =-\frac{\theta}{16r(1-r)p}+O_N(\tau p^{-2}),
\]
so its absolute value is at least $c\tau/p$ for large $p$.
The correction changes it by at most
$Cp\tau p^{-L_{\mathrm{Newt}}}=o(\tau/p)$ since $L_{\mathrm{Newt}}>2$.
It remains nonzero.
Any ball invariant under both orthogonal block groups has centre
fixed by both groups, hence centre zero. Its radial function
would be constant and would have zero monic degree-two angular
coefficient. The nonzero functional excludes such a ball.
\end{proof}

\begin{proposition}[Identification of the zero with a classical solution]
\label{inf:classical-solution}
For a real zero $(g,f)$ of \eqref{inf:nonlinear-definition} in the
small chart, set $U=1+Kg$, $\widehat\Omega=\Phi_f(B)$, and
$w=U\circ\Phi_f^{-1}$. Then $w$ is a classical solution of
\begin{equation}\label{inf:normalized-PDE}
 \Delta w+R^2w=0\quad\hbox{in }\widehat\Omega,
 \qquad w=1,\quad\nabla w=0\quad\hbox{on }\partial\widehat\Omega.
\end{equation}
It is smooth up to the boundary and real analytic there and in
the interior. For $\Omega=R\widehat\Omega$ and $u(x)=w(x/R)$,
one has $\Delta u+u=0$, boundary value one and zero boundary
gradient. The field is nonconstant.
\end{proposition}
\begin{proof}
The source $g$ is an actual $L^2$ function by
Lemma~\ref{inf:space-embedding}. Its Poisson field belongs to
$H^2(B)\cap H^1_0(B)$, so $U\in H^2(B)$ and its Dirichlet
trace is one. The completed-space trace identity
\eqref{inf:poisson-trace} makes its normal trace zero because
$g\in\Gc_p$.
The tangential trace of its gradient is the tangential derivative
of its constant Dirichlet trace, hence zero in the Sobolev trace
sense. Equivalently, finite clamped Poisson sums have zero full-gradient
trace and converge to $U$ in $H^2$; continuity of the gradient
trace gives the same conclusion on the completion. Thus the
whole gradient trace is zero.
The map $\Phi_f$ is a smooth diffeomorphism of closed domains;
composition transports $U$ to an $H^2$ field on
$\widehat\Omega$. Proposition~\ref{inf:nonlinear-map} identifies
its coefficient equation with the weak physical equation in
\eqref{inf:normalized-PDE}.
The trace chain rule
$\nabla w\circ\Phi_f=D\Phi_f^{-T}\nabla U$ transports the zero
full-gradient trace as well.

Dirichlet elliptic regularity on a smooth bounded domain applies
to $\Delta w=-R^2w$ and the constant boundary data. Starting from
$H^2$, each application improves regularity by two derivatives;
iteration gives all Sobolev orders, hence $C^\infty$ up to the
boundary in the fixed dimension; see \cite[Chs.~6, 8, 9]{inf:GT}.
The boundary is analytic by Lemma~\ref{inf:space-embedding}, the
operator has analytic constant coefficients, and its Dirichlet
data are analytic. Analytic interior regularity \cite{inf:MorreyNirenberg} and
analytic Dirichlet boundary regularity \cite{inf:MorreyBoundary}
give real analyticity up to that boundary. The scalar Laplacian
is strongly elliptic, and the Dirichlet condition is its
complementing boundary condition; the lower-order term $R^2$
has analytic coefficients. Its now classical gradient trace agrees with the zero
Sobolev trace. This proves \eqref{inf:normalized-PDE} pointwise.

The dilation gives
$\Delta_xu=R^{-2}(\Delta w)(x/R)=-u$ and preserves both boundary
conditions. The only constant solution of $\Delta u+u=0$ is zero,
which is incompatible with boundary value one. Hence $u$ is
nonconstant.
\end{proof}

\begin{remark}[The symmetry sector]\label{inf:sector-distinction}
The Dirichlet inverse throughout the construction is only in the
$O(a)\times O(b)$ invariant sector. At the exact solution, Cartesian
derivatives of $u$ satisfy the frequency-one equation and have zero
Dirichlet trace. At least one is nonzero, so the full ambient
Dirichlet operator has eigenvalue one. Since $u$ is invariant,
its $X$-derivatives
transform as the standard $O(a)$ representation and its
$Y$-derivatives as the standard $O(b)$ representation. Averaging
any of these derivatives over the corresponding reflection group
gives zero. They therefore have zero invariant projection.
Within the invariant sector the graph inverse used above includes
every mode, in particular the degree-four seed and the harmonic
source layer.
\end{remark}

\begin{proposition}[Failure of the Pompeiu property]\label{inf:Pompeiu}
Every bounded domain carrying a classical field constructed above fails
the Pompeiu property.
\end{proposition}
\begin{proof}
For any $\xi\in\R^{2p}$ with $|\xi|=1$, put
$v(x)=e^{i\xi\cdot x}$. Both $u$ and $v$ solve the frequency-one
Helmholtz equation. Green's second identity gives
\[
 0=\int_{\partial\Omega}(u\partial_\nu v-v\partial_\nu u)
   =\int_{\partial\Omega}\partial_\nu v
   =\int_\Omega\Delta v=-\int_\Omega e^{i\xi\cdot x}\,dx.
\]
Thus the Fourier transform of the indicator vanishes on the unit
sphere. Fix a unit vector $\xi_0$.
For every orthogonal map $\mathsf O$ and translation vector $x_{\mathrm{tr}}$,
\[
 \int_{x_{\mathrm{tr}}+\mathsf O\Omega}e^{i\xi_0\cdot x}\,dx
 =e^{i\xi_0\cdot x_{\mathrm{tr}}}\int_\Omega e^{i(\mathsf O^T\xi_0)\cdot y}\,dy=0.
\]
Taking real parts provides the nonzero continuous real function
$\cos(\xi_0\cdot x)$ whose integrals over all rigid copies of
$\Omega$ vanish. This is the failure of the Pompeiu property.
\end{proof}

\begin{corollary}[The bounded base construction]\label{inf:bounded-basic}
For every sufficiently large $p\in\mathcal P_{\mathrm{bd}}$, some
integer $a\in[p/2,3p/4]$ gives a bounded strictly convex analytic
non-ball domain with $O(a)\times O(2p-a)$ symmetry and the data
\eqref{inf:main-equations}. The construction uses
$0<\tau\le20$ and the cofactor order $N=4M+40$.
\end{corollary}
\begin{proof}
Lemma~\ref{inf:shifted-arithmetic} supplies unbounded inputs of both
parities; Lemma~\ref{inf:detuning} supplies their radii and
$0<\tau<20$. Theorem~\ref{inf:integer-gap} selects an integer
split for every sufficiently large member. Fix the order of
constants \eqref{inf:constant-order} and the resulting one finite
accuracy order $N$.
The centre exists at every selected split by
Theorem~\ref{inf:centres}; its gap transfers by
Proposition~\ref{inf:centre-gap-transfer}.
Theorem~\ref{inf:coefficient-inverse},
Theorem~\ref{inf:material-isomorphism}, and the second-derivative
and residual bounds give the exact zero in
Corollary~\ref{inf:Newton-zero}.
Lemmas~\ref{inf:geometry-convexity} and
\ref{inf:geometry-nonball} give its analytic strictly convex
non-ball domain with the required block symmetry.
Proposition~\ref{inf:classical-solution} gives the classical
frequency-one field, and Proposition~\ref{inf:Pompeiu} gives the
Pompeiu assertion.
There are only finitely many thresholds after the constants and
$N$ have been fixed; all are uniform in the split and in the
nonzero detuning range. Their maximum discards only finitely
many members of $\mathcal P_{\mathrm{bd}}$. Its remaining members have
unbounded distinct dimensions $2p$, as required.
\end{proof}
